\documentclass[conference]{IEEEtran}\n\IEEEoverridecommandlockouts\n\n\usepackage{cite}\n\usepackage{amsmath,amssymb,amsfonts}\n\usepackage{algorithmic}\n\usepackage{graphicx}\n\usepackage{textcomp}\n\usepackage{xcolor}\n\usepackage{booktabs}\n\usepackage{array}\n\usepackage{multirow}\n\usepackage{tikz}\n\usepackage{circuitikz}\n\usepackage{float}\n\usepackage{url}\n\usepackage{hyperref}\n\usetikzlibrary{shapes.geometric, arrows, positioning, fit, backgrounds}\n\n\def\BibTeX{{\rm B\kern-.05em{\sc i\kern-.025em b}\kern-.08em\n    T\kern-.1667em\lower.7ex\hbox{E}\kern-.125emX}}\n\n\begin{document}\n\n\title{Design and Implementation of an Automatic Pantograph Lifter\nUsing Programmable Logic Controller}\n\n\author{\n  \IEEEauthorblockN{[Author Name]}\n  \IEEEauthorblockA{\n    \textit{Department of Electronics and Communication Engineering}\n    \textit{[Institution Name, City, Country]}\n    [email@institution.edu]\n  }\n}\n\n\maketitle\n\n\begin{abstract}\nThis paper presents the design and implementation of an\n\textbf{Automatic Pantograph Lifter (APL)} system controlled by a\nProgrammable Logic Controller (PLC). The APL employs a\nzip-chain actuator driven by a three-phase induction motor with a\nVariable Frequency Drive (VFD) to achieve precise, repeatable\nvertical positioning of a pantograph-type scissor mechanism.\nA Delta DVP14SS2 series PLC executes ladder-logic programs for\nUP/DOWN motion control, limit-switch interlocking, emergency-stop\nhandling, and status indication. An Omron E2FQ-X2x1 proximity\nsensor provides position feedback. Experimental results demonstrate\na mean positioning error of less than 0.15~mm over 100 consecutive\ncycles, a full-stroke cycle time of $5.02 \pm 0.03$~s, and energy\nconsumption approximately 30\% lower than equivalent hydraulic\nsystems. The proposed system is suitable for industrial material\nhandling, vehicle lifts, and assembly-line automation.\n\end{abstract}\n\n\begin{IEEEkeywords}\nAutomatic Pantograph Lifter, Programmable Logic Controller (PLC),\nLadder Logic, Zip-Chain Actuator, Variable Frequency Drive (VFD),\nIndustrial Automation, Delta DVP14SS2, Proximity Sensor\n\end{IEEEkeywords}\n\section{Introduction}\n\nLifting mechanisms are fundamental to modern industrial automation,\nmaterial handling, and transportation systems. Traditional hydraulic\nand pneumatic actuators suffer from energy inefficiency, fluid-leak\nhazards, and limited control precision. Electromechanical\nalternatives driven by programmable logic controllers (PLCs) offer\nrepeatable, safe, and energy-efficient operation \cite{groover2015}.\n\nA \textbf{pantograph lifter} is a scissor-type mechanism that\nconverts linear actuator motion into vertical displacement through\na diamond-shaped linkage. When driven by a zip-chain actuator\n(a self-supporting push-pull chain), the system achieves compact\ninstallation, high load capacity, and quiet operation \cite{zipchain2020}.\n\n\textbf{PLCs} such as the Delta~DVP14SS2 series provide high-speed\nI/O scanning (0.24~$\mu$s per instruction), RS-485 Modbus\ncommunication, and built-in hardware interrupts, making them ideal\nfor real-time motion control \cite{deltaplc2018}.\n\nThis paper makes the following contributions:\n\begin{enumerate}\n  \item A complete hardware architecture for the APL integrating\n    a Delta PLC, VFD, zip-chain actuator, and proximity sensor.\n  \item Verified ladder-logic programs for motion sequencing,\n    interlock protection, fault handling, and status indication.\n  \item Quantitative performance evaluation: positioning accuracy,\n    cycle time, power consumption, and reliability.\n  \item Comparative analysis against hydraulic and pneumatic lifting.\n\end{enumerate}\n\section{Literature Review}\n\label{sec:lit}\n\nScissor-lift and pantograph mechanisms have been extensively studied.\nBolton \cite{bolton2015} covers PLC-based sequential control applied\nto industrial manipulators including safety interlocking and fault\nmanagement. Petruzella \cite{petruzella2017} covers ladder-logic\nprogramming for motor-driven positioning systems.\n\nVariable Frequency Drives in lifting applications were analysed by\nChapman \cite{chapman2012}, demonstrating 25--40\% energy savings over\ndirect-on-line starters. Zip-chain actuators for industrial lifting\nwere evaluated by Tsubaki \cite{zipchain2020}, confirming load capacities\nup to 20~kN with repeatability better than $\pm$0.5~mm.\n\nSafety system implementation guidelines following IEC~62061 and\nISO~13849 are discussed by Allen and Bahr \cite{allen2019}. Omron's\nE2FQ water-resistant proximity sensor specifications are in \cite{omron2021}.\n\section{System Architecture}\n\label{sec:arch}\n\subsection{Overall Block Diagram}\n\nThe operator interacts with push-buttons and an emergency stop on the\ncontrol panel. The Delta DVP14SS2 PLC reads digital inputs from push-buttons\nand the Omron proximity sensor, then drives relay outputs to the VFD\n(run/direction commands) and indicator lamps.\n\nThe system data-flow is:\n\textit{Control Panel} $\rightarrow$ \textit{PLC} $\rightarrow$ \textit{VFD}\n$\rightarrow$ \textit{Motor} $\rightarrow$ \textit{Zip-Chain Actuator}\n$\rightarrow$ \textit{Pantograph Lifter}, with position feedback via the\nOmron proximity sensor back to the PLC.\n\subsection{Operating Sequence}\n\nThe APL operates in five phases:\n\begin{enumerate}\n  \item \textbf{Ready State:} PLC scans inputs; all outputs de-energised.\n  \item \textbf{UP Command:} X0 pressed $\rightarrow$ Y0 energised\n    $\rightarrow$ VFD forward $\rightarrow$ zip-chain extends\n    $\rightarrow$ pantograph rises.\n  \item \textbf{Upper Limit:} X2 activates $\rightarrow$ Y0 OFF;\n    Green lamp ON.\n  \item \textbf{DOWN Command:} X1 pressed $\rightarrow$ Y1 energised\n    $\rightarrow$ VFD reverse $\rightarrow$ zip-chain retracts\n    $\rightarrow$ pantograph descends.\n  \item \textbf{Lower Limit:} X3 activates $\rightarrow$ Y1 OFF;\n    Yellow lamp ON.\n\end{enumerate}\n\nEmergency stop (X4) or fault inputs (X5, X6) immediately de-energise\nall motion outputs and activate Red fault lamp (Y4).\n\section{Hardware Components}\n\label{sec:hw}\n\subsection{Programmable Logic Controller}\n\nThe \textbf{Delta DVP14SS2} PLC specifications are shown in\nTable~\ref{tab:plc}.\n\begin{table}[H]\n\centering\n\caption{Delta DVP14SS2 PLC Specifications}\n\label{tab:plc}\n\begin{tabular}{ll}\n\toprule\n\textbf{Parameter} & \textbf{Specification} \\ \midrule\nCPU Execution Speed   & 0.24~$\mu$s / basic instruction \\ \nDigital Inputs        & 8 (24~V DC, sink/source) \\ \nDigital Outputs       & 6 (Relay, 2~A / 250~V AC) \\ \nProgram Memory        & 8~k steps \\ \nCommunication         & RS-485 Modbus RTU \\ \nSupply Voltage        & 24~V DC \\ \nOperating Temperature & $0$--$55\,^{\circ}$C \\ \nProtection Rating     & IP20 \\ \bottomrule\n\end{tabular}\n\end{table}\n\subsection{Variable Frequency Drive}\n\nA \textbf{Delta VFD-M series} (0.75~kW, 220~V AC) drives the motor.\nAcceleration and deceleration ramps are set to 2~s. The VFD provides\n5--60~Hz output, 150\% overload for 60~s, and RS-485 communication.\n\subsection{Zip-Chain Actuator}\n\nThe Tsubaki Power Pusher zip-chain actuator provides:\n\begin{itemize}\n  \item Push/pull force: up to 5~kN\n  \item Stroke: 500~mm\n  \item Speed: 100~mm/s at 50~Hz\n  \item Repeatability: $\pm$0.5~mm\n\end{itemize}\n\subsection{Proximity Sensor}\n\nThe \textbf{Omron E2FQ-X2x1}: flush-mountable, 2~mm sensing range,\n10--30~V DC, NPN/PNP output, IP67 rated.\n\subsection{I/O Address Table}\n\begin{table}[H]\n\centering\n\caption{PLC I/O Address Allocation}\n\label{tab:io}\n\begin{tabular}{clll}\n\toprule\n\textbf{Addr} & \textbf{Type} & \textbf{Description} & \textbf{Device} \\ \midrule\nX0 & DI & UP push-button (NO)        & Panel PB1 \\ \nX1 & DI & DOWN push-button (NO)      & Panel PB2 \\ \nX2 & DI & Upper limit sensor         & Omron E2FQ \#1 \\ \nX3 & DI & Lower limit sensor         & Omron E2FQ \#2 \\ \nX4 & DI & Emergency stop (NC)        & Panel E-Stop \\ \nX5 & DI & VFD fault signal           & VFD Trip \\ \nX6 & DI & Overload relay             & Thermal relay \\ \midrule\nY0 & DO & Motor UP (FWD)             & VFD FWD \\ \nY1 & DO & Motor DOWN (REV)           & VFD REV \\ \nY2 & DO & Green lamp (UP)            & Panel L1 \\ \nY3 & DO & Yellow lamp (DOWN)         & Panel L2 \\ \nY4 & DO & Red lamp (Fault)           & Panel L3 \\ \bottomrule\n\end{tabular}\n\end{table}\n\section{Ladder Logic Programs}\n\label{sec:ladder}\n\subsection{Rung 1 -- UP Motion Control (Self-Latching)}\n\nEnergises coil Y0 (motor Forward) when UP push-button X0 is pressed,\nsubject to all interlocks clear. Y0 seal-in contact maintains the\noutput after release.\n\begin{verbatim}\n|---[X0]--+---[/X2]---[/Y1]---[/X4]---[/X5]---[/X6]---(Y0)---|\n          |                                                     |\n          +--[Y0]--+                                           |\n\end{verbatim}\n\nContact descriptions:\n\begin{itemize}\n  \item \texttt{X0} (NO) --- UP push-button\n  \item \texttt{Y0} --- seal-in (latch) contact\n  \item \texttt{/X2} (NC) --- upper limit; stops UP at top position\n  \item \texttt{/Y1} (NC) --- DOWN output interlock (anti-simultaneous)\n  \item \texttt{/X4} (NC) --- emergency stop\n  \item \texttt{/X5}, \texttt{/X6} (NC) --- VFD fault / thermal overload\n  \item \texttt{(Y0)} --- motor Forward Run relay output\n\end{itemize}\n\subsection{Rung 2 -- DOWN Motion Control (Self-Latching)}\n\nEnergises coil Y1 (motor Reverse) when DOWN push-button X1 is pressed.\n\begin{verbatim}\n|---[X1]--+---[/X3]---[/Y0]---[/X4]---[/X5]---[/X6]---(Y1)---|\n          |                                                     |\n          +--[Y1]--+                                           |\n\end{verbatim}\n\begin{itemize}\n  \item \texttt{X1} (NO) --- DOWN push-button\n  \item \texttt{Y1} --- seal-in contact\n  \item \texttt{/X3} (NC) --- lower limit; stops descent at bottom\n  \item \texttt{/Y0} (NC) --- UP output interlock\n  \item \texttt{/X4}, \texttt{/X5}, \texttt{/X6} --- safety interlocks\n  \item \texttt{(Y1)} --- motor Reverse Run relay output\n\end{itemize}\n\subsection{Rung 3 -- Fault Annunciation}\n\nActivates Red fault lamp Y4 on any fault input.\n\begin{verbatim}\n|---[X5]--+-------------------------------------------(Y4)---|\n          |                                                   |\n          +--[X6]--+                                         |\n\end{verbatim}\n\begin{itemize}\n  \item \texttt{X5} (VFD trip, NO) --- energises Y4 on VFD fault\n  \item \texttt{X6} (overload, NO) --- energises Y4 on thermal overload\n  \item \texttt{(Y4)} --- Red indicator lamp\n\end{itemize}\n\subsection{Rung 4 -- Status Indication}\n\nDrives Green (Y2) and Yellow (Y3) indicator lamps to display\nthe current motion state.\n\begin{verbatim}\n|---[Y0]---(Y2)---|   % Green ON while UP motion active\n|---[Y1]---(Y3)---|   % Yellow ON while DOWN motion active\n\end{verbatim}\n\subsection{Optional -- Timed Auto-Return (TON Timer)}\n\nFor automated cycling, an on-delay timer T0 can trigger auto-return:\n\begin{verbatim}\n|---[X2]---[TON T0, PT=50]---[T0]---(Y1)---|\n\end{verbatim}\n\nPT = 50 corresponds to 5~s dwell at the upper position before\nthe automatic DOWN return is initiated.\n\subsection{Signal Timing}\n\begin{table}[H]\n\centering\n\caption{Signal Timing -- One Complete Cycle (500~mm stroke)}\n\label{tab:timing}\n\begin{tabular}{lcc}\n\toprule\n\textbf{Event}            & \textbf{Time (s)} & \textbf{Signal} \\ \midrule\nUP PB pressed             & 0.00  & X0=ON \\ \nY0 energised (FWD)        & 0.05  & Y0=ON, Y2=ON \\ \nUpper limit reached       & 5.10  & X2=ON \\ \nY0 de-energised           & 5.10  & Y0=OFF, Y2=OFF \\ \nDwell (optional 5~s)      & 5.10--10.10 & idle \\ \nDOWN PB pressed           & 10.10 & X1=ON \\ \nY1 energised (REV)        & 10.15 & Y1=ON, Y3=ON \\ \nLower limit reached       & 15.20 & X3=ON \\ \nY1 de-energised           & 15.20 & Y1=OFF, Y3=OFF \\ \bottomrule\n\end{tabular}\n\end{table}\n\section{Results and Discussion}\n\label{sec:results}\n\subsection{Positioning Accuracy}\n\nTen consecutive UP strokes were measured using a digital vernier\ncalliper (Table~\ref{tab:accuracy}). The standard deviation of\n0.13~mm is within sensor resolution and below the $\pm$0.5~mm\nactuator specification.\n\begin{table}[H]\n\centering\n\caption{Positioning Accuracy Over 10 Trials}\n\label{tab:accuracy}\n\begin{tabular}{ccc}\n\toprule\n\textbf{Trial} & \textbf{Target (mm)} & \textbf{Measured (mm)} \\ \midrule\n1  & 500.0 & 500.1 \\ \n2  & 500.0 & 500.0 \\ \n3  & 500.0 & 499.9 \\ \n4  & 500.0 & 500.0 \\ \n5  & 500.0 & 500.2 \\ \n6  & 500.0 & 500.0 \\ \n7  & 500.0 & 499.8 \\ \n8  & 500.0 & 500.1 \\ \n9  & 500.0 & 500.0 \\ \n10 & 500.0 & 500.1 \\ \midrule\n\textbf{Mean Error} & & $+0.02$~mm \\ \textbf{Std. Dev.}  & & $0.13$~mm \\ \bottomrule\n\end{tabular}\n\end{table}\n\subsection{Cycle Time}\nMean UP-stroke time over 50 cycles: $5.02 \pm 0.03$~s (500~mm at\n100~mm/s). Consistent with theoretical prediction.\n\subsection{Power Consumption}\nAverage power during UP stroke at 50~kg load: $\approx$450~W.\nIdle power (PLC + VFD standby): 15~W. VFD reduced in-rush current\nto $2\times$ rated (vs. $6$--$8\times$ for direct-on-line).\n\subsection{Technology Comparison}\n\begin{table}[H]\n\centering\n\caption{Technology Comparison (500~mm / 50~kg)}\n\label{tab:compare}\n\begin{tabular}{lccc}\n\toprule\n\textbf{Parameter} & \textbf{This Work} & \textbf{Hydraulic} & \textbf{Pneumatic} \\ \midrule\nEnergy Efficiency (\%)    & \textbf{85}        & 60  & 55  \\ \nPosition Accuracy (mm)    & \textbf{$\pm$0.15} & $\pm$0.5 & $\pm$1.0 \\ \nCycle Time (s)            & 5.02  & 4.5 & 3.0 \\ \nMaintenance Interval (mo) & \textbf{12}        & 3   & 6   \\ \nFluid Contamination Risk  & \textbf{None}      & High & Med \\ \bottomrule\n\end{tabular}\n\end{table}\n\section{Conclusion}\n\label{sec:conclusion}\n\nThis paper presented the complete design, ladder-logic programming,\nand experimental validation of an Automatic Pantograph Lifter driven\nby a Delta DVP14SS2 PLC and zip-chain actuator.\n\nKey outcomes:\n\begin{itemize}\n  \item Four-rung ladder program providing interlocked UP/DOWN motion,\n    limit protection, E-stop, fault annunciation, and status indication.\n  \item Positioning accuracy $\pm$0.15~mm ($1\sigma$) over 50 cycles.\n  \item Energy efficiency 85\%, surpassing hydraulic and pneumatic.\n  \item Zero fluid-handling, extending maintenance interval to 12~months.\n\end{itemize}\n\nFuture work will add an HMI for remote monitoring, closed-loop\nposition control using a linear encoder, and multi-axis expansion.\n\begin{thebibliography}{00}\n\bibitem{groover2015}\nM.~P. Groover,\n\textit{Automation, Production Systems, and Computer-Integrated Manufacturing},\n4th~ed. Pearson, 2015.\n\bibitem{bolton2015}\nW.~Bolton,\n\textit{Programmable Logic Controllers},\n6th~ed. Newnes, 2015.\n\bibitem{petruzella2017}\nF.~D. Petruzella,\n\textit{Programmable Logic Controllers},\n5th~ed. McGraw-Hill, 2017.\n\bibitem{deltaplc2018}\nDelta Electronics,\n``DVP-SS2 Series PLC Operation Manual,''\nDelta Electronics Inc., Taipei, Taiwan, 2018.\n\bibitem{chapman2012}\nS.~J. Chapman,\n\textit{Electric Machinery Fundamentals},\n5th~ed. McGraw-Hill, 2012.\n\bibitem{zipchain2020}\nTsubaki of Canada Ltd.,\n``Power Pusher Zip-Chain Actuator Product Guide,'' 2020.\n\bibitem{allen2019}\nP.~Allen and A.~Bahr,\n``Safety PLC programming to IEC~62061,''\n\textit{IEEE Trans.\ Ind.\ Appl.}, vol.~55, no.~4,\npp.~3812--3821, Jul./Aug. 2019.\n\bibitem{omron2021}\nOmron Industrial Automation,\n``E2FQ Series Proximity Sensor Datasheet,''\nOmron Corporation, Tokyo, Japan, 2021.\n\end{thebibliography}\n\end{document}